\section{Discussion}\label{s:discussion}
In Section~\ref{sec:accuracy}, we could observe clear differences between the accuracy of different simulators and power models.
Although \textit{OpenDC} and \textit{DISSECT-CF} use different values as the constant power consumption value, both their constant models performed quite well. For their constant values, \textit{OpenDC} a user-defined value while \textit{DISSECT-CF} uses the idle power consumption value. A reason for the good accuracy of the constant models performing well could be the fact that the power usage values given to these models were derived from the superset that contains the same traces they then simulated. This suggests that these models are good at modelling stable, predictable behaviour, but perhaps this accuracy does not carry over to more unpredictable system or applications.
Between the \textit{Sock Shop} and \textit{Train Ticket} traces, we can see that with the higher complexity of \textit{Train Ticket}, simulators produce more similar results, also across power models. \textit{CloudSim}, \textit{NetworkCloudSim} and \textit{DISSECT-CF} with the linear power model drastically overestimate the power usage, while \textit{DISSECT-CF} with the constant power model, \textit{SimGrid} and \textit{OpenDC} with the non-constant power models underestimate the power usage. An explanation for this could be that \textit{Train Ticket} has more services than \textit{Sock Shop} in combination with the CPU usage value being given per service. If there is a per-service margin of error and a larger number of services, the total error increases with a larger number of service. As microservice systems in industry can have many more services than \textit{Train Ticket}, and it is precisely these systems that consume the largest amounts of energy, this indicates a shortcoming of the simulators towards simulating microservice systems and one of the areas for improvement in the future regarding cloud simulators.

\subsection{Future Work}
\begin{itemize}
    \item \textbf{Simulators calibration}: \rone{as we will discuss in Section~\ref{sec:internal_threats}, for all simulators we configured the hardware-dependent idle and maximum power consumption values of the host machines, while we kept all the other simulation parameters to their default values.  An important direction for future research is the systematic investigation of the impact of simulator calibration on energy-estimation accuracy and how different calibration strategies affect the predictive performance of different simulators across diverse workloads, infrastructures, and/or deployment scenarios. Building on such evidence, it would be possible to develop a decision model that supports researchers and practitioners in (i) selecting the most suitable cloud simulator for a given energy-evaluation objective and also (ii) identifying appropriate calibration strategies for that simulator. Such a decision model could provide guidelines on the expected trade-offs among calibration effort, simulation complexity, reproducibility, and estimation accuracy, ultimately leading to more accurate energy evaluations of microservice systems.}
    \item \textbf{Dynamic Power models}: as shown in \autoref{sec:resultsEnergy}, the power model that best describes the energy consumption of a microservice system can change over time. Therefore, research into dynamic power models that change depending on the context could be an interesting area of improvement for energy simulation of microservice systems.
    \item \textbf{Service-level energy simulation}: energy simulation in cloud simulators is currently only available on a host-level. This makes the use of these simulators for microservice research limited, making it very hard to carry out root cause analysis into extraordinarily high energy usage levels caused by certain parts of the microservice system, for example. \rone{Assessing service-level energy-consumption estimation remains an important future research direction for researchers working on cloud simulators and would require suitable fine-grained power models for individual services and publicly available datasets providing reliable per-service energy measurements (the ground truth dataset we constructed for this study is a first step towards this direction).}
    \item \textbf{Scalability}: the simulators studied have scalability issues both in terms of accuracy decreasing with increased system complexity and simulation time increasing when communication is simulated. The limited scalability of cloud simulators currently available has been noted by others (e.g. \cite{martone_evaluation_2023}) and is particularly relevant to the microservices context, as microservice systems have a lot of communication and are often complex systems comprised of many services. Furthermore, it is especially the big, complex microservice systems for which the energy usage should be reduced.
    \item \textbf{Increased Usability}: often, cloud simulators are a tool used for achieving a certain goal instead of the goal themselves. Out of the 5 simulators discussed here, none have a GUI usable in a research context, and for all except \textit{OpenDC}, users have to largely code the experiment to be simulated themselves. Although documentation or examples are available, this still means that simulation takes considerable effort and users have to be able to code in the respective simulator's language. The threshold for using simulators would be decreased if they were equipped with a usable GUI. Furthermore, the ability to load traces from file could be implemented for all simulators, as well as some basic energy usage reporting such that the users do not have to implement these themselves. 
\end{itemize}